%=== APPENDIX ===

\chapter*{Appendix A: COVID-19 Scenario Data}
\addcontentsline{toc}{chapter}{Appendix A: COVID-19 Scenario Data}


%Mapping of Sectors

\begin{table}[H]
\centering
\small
\setlength{\tabcolsep}{3pt}
\renewcommand{\arraystretch}{1.1}
\begin{tabular}{@{}ccc@{}}
\toprule
\textbf{Sector ID} & \textbf{PC1} & \textbf{PC2} \\ \midrule
3 & 0.772976374 & -0.158540904 \\
2 & 0.488575487 & 0.547564001 \\
5 & 0.230631604 & 0.602425551 \\
8 & 0.173952844 & 0.31984955 \\
4 & 0.16376151 & 0.197548491 \\
10 & 0.141148665 & 0.185477914 \\
7 & 0.107051197 & 0.232037887 \\
11 & 0.100445788 & 0.154054165 \\
1 & 0.082357584 & 0.203204452 \\
9 & 0.065932369 & 0.120308858 \\
12 & 0.02097817 & 0.033819003 \\
6 & 0.01904942 & 0.04009044 \\
15 & 0.010160651 & 0.015616709 \\
14 & 0.004471868 & 0.006256627 \\
13 & 0.002201633 & 0.00402146 \\ \bottomrule
\end{tabular}
\caption{Sector PCA loadings (COVID-19, 15 sectors)}
\label{tab:sector-mapping}
\end{table}



%Cross-section of technical coefficient matrix

\begin{table}[H]
\centering
\small
\setlength{\tabcolsep}{3pt}
\renewcommand{\arraystretch}{1.1}
\begin{tabular}{@{}
>{\columncolor[HTML]{FFFFFF}}c 
>{\columncolor[HTML]{FFFFFF}}c 
>{\columncolor[HTML]{FFFFFF}}c 
>{\columncolor[HTML]{FFFFFF}}c @{}}
\toprule
{\color[HTML]{000000} } & {\color[HTML]{000000} \begin{tabular}[c]{@{}c@{}}Agriculture, Fishing, Quarrying, \\ Utilities and Sewerage \\ \& Waste Management\end{tabular}} & {\color[HTML]{000000} Manufacturing} & {\color[HTML]{000000} Construction} \\ \midrule
{\color[HTML]{000000} \begin{tabular}[c]{@{}c@{}}Agriculture, Fishing, Quarrying, Utilities \\ and Sewerage \& Waste Management\end{tabular}} & {\color[HTML]{000000} 4,308.828700} & {\color[HTML]{000000} 2,985.482100} & {\color[HTML]{000000} 198.954100} \\
{\color[HTML]{000000} Manufacturing} & {\color[HTML]{000000} 498.915600} & {\color[HTML]{000000} 59,777.284000} & {\color[HTML]{000000} 4,939.792300} \\
{\color[HTML]{000000} Construction} & {\color[HTML]{000000} 166.973900} & {\color[HTML]{000000} 498.152500} & {\color[HTML]{000000} 23,583.027800} \\ \bottomrule
\end{tabular}
\caption{Cross-section of technical coefficient matrix (COVID-19)}
\label{tab:tech-coeff-cross-section}
\end{table}

%Initial inoperability

\begin{table}[H]
\centering
\small
\setlength{\tabcolsep}{3pt}
\renewcommand{\arraystretch}{1.1}
\begin{tabular}{@{}ccc@{}}
\toprule
\rowcolor[HTML]{FFFFFF} 
\textbf{Sector ID} & \textbf{Sector} & \textbf{Initial inoperability} \\ \midrule
1 & Agriculture, Fishing,   Quarrying, Utilities and Sewerage \& Waste Management & 0.020592 \\
2 & Manufacturing & 0.033293 \\
3 & Construction & 0.155924 \\
4 & Wholesale  \&    Retail  Trade & 0.061522 \\
5 & Transportation  \&    Storage & 0.032185 \\
6 & Accommodation  \&    Food  Services & 0.184162 \\
7 & Information  \&    Communications & 0.006077 \\
8 & Financial  \&    Insurance  Services & 0.009252 \\
9 & Real  Estate    Services & 0.026593 \\
10 & Professional  Services & 0.105436 \\
11 & Administrative  \&    Support  Services & 0.123510 \\
12 & Public  Administration  \&    Education & 0.115083 \\
13 & Health  \&    Social  Services & 0.100575 \\
14 & Arts,  Entertainment  \&    Recreation & 0.177582 \\
15 & Other  Community,    Social  \&  Personal    Services & 0.439401 \\ \bottomrule
\end{tabular}
\caption{Sector Initial Inoperability (COVID-19)}
\label{tab:sector-inoperability}
\end{table}


\chapter*{Appendix B: Manpower Scenario Data}
\addcontentsline{toc}{chapter}{Appendix B: Manpower Scenario Data}

The manpower disruption scenario uses the 2022 Singapore IO table with 15 sectors. Due to the large number of sectors, only summary statistics are presented here.

\begin{table}[H]
\centering
\small
\begin{tabular}{@{}lcc@{}}
\toprule
\textbf{Parameter} & \textbf{COVID-19} & \textbf{Manpower} \\
\midrule
Number of sectors ($n$) & 15 & 15 \\
IO table year & 2019 & 2022 \\
Mean initial inoperability & 0.10608 & varies \\
Max initial inoperability & 0.43940 & varies \\
Demand shock ($c^*$) & $> 0$ & $= 0$ \\
Lockdown duration & 55 days & 55 days \\
Total horizon & 751 days & 751 days \\
PCA $\times\, x_i$ top-5 sectors & 2, 5, 4, 8, 3 & 2, 5, 4, 7, 8 \\
DIIM top-5 sectors & by simulation & by simulation \\
\bottomrule
\end{tabular}
\caption{Summary of scenario parameters}
\label{tab:scenario-summary}
\end{table}


\chapter*{Appendix C: Code}
\addcontentsline{toc}{chapter}{Appendix C: Code}

The R code for all simulations and analyses is available in the project repository. Key scripts include:

\begin{itemize}
    \item \texttt{functions.R}: Core DIIM function, data loading, and simplified method implementations (Total Output, PCA $\times\, x_i$, PageRank $\times\, x_i$, BL $\times\, x_i$, FL $\times\, x_i$)
    \item \texttt{pca.R}: PCA analysis on the integrated IO matrix
    \item \texttt{simulations/decision\_rules.R}: Monte Carlo framework, logistic regression, and decision boundary rules
    \item \texttt{simulations/enhanced\_comparison.R}: Multi-PC PCA, varying top-$k$ sweep, and baseline vs structured method comparison
    \item \texttt{simulations/run\_simulations.R}: Parameter sweeps for severity ($c^*$) and lockdown duration impacts
    \item \texttt{manpower-disruption/manpower-disruption.R}: Manpower scenario DIIM simulation implementation
\end{itemize}


%=== END OF APPENDIX ===
\newpage
